%% ============================================================================
\section{Cross-Coin Evidence}
\label{sec:placebo}

Reserve exposure orders the three coins ex ante: USDC held direct SVB exposure, DAI was
pinned to USDC through the peg-stability module, and USDT held none. That gives
the ex-ante fragility ordering USDC, DAI, USDT. Troughs are \$0.877, \$0.886, and
\$0.992 in that order. The deviations from par are $1233$, $1141$, and $76$~bps.
DAI's mechanical link to USDC leaves one free
comparison, so the realized ordering has probability $1/2$ and illustrates the ex-ante
ordering without testing it. USDC and DAI clear the $1\%$
break threshold, a deviation of at least one cent from par, by factors of
$12$ and $11$. USDT does not clear it, so the placebo criterion is satisfied: the
exposed coins break and the unexposed coin does not.

\noindent\textbf{USDT as the unexposed control.} USDT was bid to a premium while USDC and
DAI fell. That pattern is consistent with the crypto safe-haven behavior documented for
USDT~\cite{Baur_2021}. The premium is venue- and time-specific. On Coinbase
\texttt{USDT-USD}~\cite{coinbase2026exchangeapi} it peaked at $+266$~bps at 02:00~UTC on
11~March, 71 minutes before
Circle's 03:11~UTC disclosure. Over the window \$1.24B of \$1.52B traded above par, and
that peak is one hour of it. The composite feed peaks later and higher, at $+270$~bps at
19:00~UTC on 13~March. That peak falls after the joint federal backstop of the previous
evening, so it is post-backstop rather than intra-run. The
Coinbase premium responds to the SVB failure itself, because Circle's
USDC-specific disclosure had not yet been made. A
common-factor repricing of stablecoin risk would move all three coins down
together, so it cannot produce that premium. Rational substitution out of a gated
claim can, and the two are not separated here.
